\section{从手工特征孤岛到统一架构}

上一章说明了本讲要回答“为什么一门 Transformer 课会存在”。本章给出历史答案：十多年前，各领域不仅使用不同数据，还使用不同特征、词汇和软件栈；深度学习先统一了训练范式，Transformer 又进一步统一了主干架构。统一降低了跨领域学习成本，也让算力、数据与规模成为共同变量。

\subsection{2011：把特征工程拼成一个脆弱 pipeline}

Karpathy 从自己进入计算机视觉的经历讲起。问题不是研究者“不聪明”，而是系统把表示设计交给大量人工 descriptor：每篇论文贡献一个局部特征，实践者把它们全部提取、拼接，再在顶部训练分类器。这种流程难以联合优化，也难以判断错误来自特征、数据还是分类器。

\lecturefigure{slide-11-historical-context.jpg}{历史背景：从 2011 年的视觉工作流回看 2023 年。}{Stanford Online 官方视频 00:10:39--00:11:17。}

\teachervoice{他用《指环王》中回忆往昔的 Bilbo 自嘲，并问学生是否意识到 2023 年进入 AI 有多幸运。幽默背后的判断是：今天习以为常的统一工具链并非自然存在，而是表示学习、规模化训练和共享架构逐步形成的结果。}

\lecturefigure{slide-12-scene-2011.jpg}{“Scene: it is 2011”：手工视觉 pipeline 的时代入口。}{Stanford Online 官方视频 00:11:17--00:11:30。}

\term{handcrafted feature}（手工特征）是由研究者预先规定的图像统计量，如局部梯度、纹理、颜色直方图或几何模板；\term{descriptor} 是把局部区域编码成向量的规则。它们包含强 inductive bias，却不能像端到端神经网络那样让所有阶段围绕最终损失共同调整。

\lecturefigure{slide-13-old-computer-vision-paper.jpg}{旧式视觉论文中的多页特征与工程细节。}{Stanford Online 官方视频 00:11:30--00:12:11。}

\begin{knowledgebox}{读图：不要把密集术语误当成“模型更深”}
页面密度主要来自 feature zoo：不同 descriptor 各自提取边缘、纹理、颜色或空间关系，再把向量送入 SVM。\term{SVM}（Support Vector Machine，支持向量机）是在特征空间寻找分类边界的监督模型。它可以很强，但上游表示通常固定，整个 pipeline 不是一个由同一目标端到端优化的可微系统。
\end{knowledgebox}

Karpathy 的尖锐评价是：代码从四处收集，系统不仅复杂，而且会产生今天看来像 bug 的荒谬预测。当时实践者有时只把它当作不可避免的偶发失败。这个对比体现了工程文化的变化：当模型、数据和训练路径更统一后，错误更可能被定位、复现与系统改进。

\lecturefigure{slide-14-old-lstm-paper.jpg}{不同领域的旧论文：序列模型有自己的术语、图示与实现传统。}{Stanford Online 官方视频 00:12:23--00:12:56。}

LSTM 论文与视觉论文不仅处理不同对象，也使用不同问题分解方式。视觉研究者看到 hidden state、gates、time unrolling 时，要重新学习整套语言；序列研究者阅读视觉 descriptor 与几何 pipeline 也面临同样障碍。跨领域迁移的成本因此不只来自数学难度，还来自概念接口不统一。

\lecturefigure{slide-15-old-language-model-paper.jpg}{旧式语言模型论文：另一套任务、符号与评价传统。}{Stanford Online 官方视频 00:12:23--00:12:56。}

语言模型的核心目标其实可以非常统一：根据上下文预测下一个 token。但在早期论文中，n-gram、平滑、词形、句法和任务特定特征常形成独立技术栈。深度学习的突破之一，是让 embedding 与预测器一起从数据学习，使“上下文如何表示”不再完全依赖人工枚举。

\lecturefigure{slide-16-old-neural-architecture.jpg}{旧式神经架构图：各领域网络仍呈现显著不同的结构语言。}{Stanford Online 官方视频 00:12:23--00:12:56。}

这张图的教学重点不是逐个辨认方框，而是观察信息路径。早期系统常把任务假设直接硬编码进 topology：图像用局部二维连接，文本用时间递归，语音用专门前端。这样的 bias 可能提高样本效率，却让新模态、新传感器或长距离关系需要重新设计骨架。

\lecturefigure{slide-17-ai-areas-different.jpg}{2012 年前：视觉、语言、语音、翻译与 RL 像不同学科。}{Stanford Online 官方视频 00:12:56--00:13:17。}

\begin{warningbox}{统一不等于抹掉领域知识}
共同的神经网络语言降低了阅读门槛，但不能推出所有领域只需同一数据处理。图像的二维邻域、语音的时间频率结构、RL 的行为条件与蛋白质的几何约束仍需进入 tokenization、position/type embedding、mask、objective 或 specialized module。统一的是可组合主干，不是现实世界本身。
\end{warningbox}

\subsection{2012：规模化学习统一训练范式}

前面的问题是每个领域都在手工定义表示。ImageNet 时代展示了另一条路线：大数据集、大模型和加速硬件可以让神经网络自己学习层级表示。研究重心因此从“再发明一个 descriptor”转向“怎样获得数据、计算和稳定优化”，这一 recipe 随后复制到语音、翻译与 RL。

\lecturefigure{slide-18-imagenet-2012.jpg}{ImageNet 2012：规模、数据与 GPU 训练改变视觉研究重心。}{Stanford Online 官方视频 00:12:56--00:13:33。}

\begin{knowledgebox}{读图：真正改变的是优化闭环}
输入像素经过共享参数的多层网络得到预测，分类损失再通过 backpropagation 调整全部层。表示不再是冻结的人工接口，而是与任务目标共同学习。这里的关键不只是卷积本身，而是“可微模型 + 大规模数据 + 并行计算 + 端到端损失”形成可复制闭环。
\end{knowledgebox}

\lecturefigure{slide-19-areas-read-more-similar.jpg}{深度学习之后：不同领域论文开始共享参数、优化器与训练语言。}{Stanford Online 官方视频 00:13:33--00:14:10。}

Karpathy 强调，跨领域论文突然出现可识别的共同词汇：parameters、optimizer、loss、neural network。统一表示框架产生知识复用：视觉研究者可以理解机器翻译的优化问题，NLP 研究者也可以借鉴残差连接和规模化训练。共同工具链使领域边界变薄，为 2017 年之后的架构统一创造条件。

\subsection{2017：连主干架构也开始收敛}

深度学习先统一“怎样训练”，Transformer 再统一“训练什么样的主干”。Karpathy 的一阶近似是：把同一个架构 copy-paste 到不同领域，主要改变数据的 chunking 与输入方式。这个说法故意夸张，但准确指出 attention 允许模型在一个元素集合上学习动态关系。

\lecturefigure{slide-20-transformer-2017.jpg}{2017 年 Transformer：一篇机器翻译论文提出可复用架构包。}{Stanford Online 官方视频 00:13:57--00:14:21。}

\begin{knowledgebox}{读图：原论文图不只是 attention}
左侧 encoder 与右侧 decoder 都包含 embedding、position、multi-head attention、feed-forward、residual add 与 normalization；decoder 还需要 causal mask 和 cross-attention。标题突出 attention，是因为它替代 recurrence 成为主要通信路径；模型能训练与扩展则依赖整套 package。
\end{knowledgebox}

\lecturefigure{slide-21-transformer-across-fields.jpg}{同一 Transformer 图在视觉、语言、语音、RL 与图学习中反复出现。}{Stanford Online 官方视频 00:14:17--00:14:40。}

这页应按“接口层—共享层—任务层”阅读：接口层把 patch、word/subword、spectrogram slice、state/action/reward 或 graph element 映射成向量；共享层用 attention 和 MLP 更新向量；任务层把输出解释为分类、生成、控制或结构预测。跨领域复用发生在中间，输入与目标仍然决定模型学到什么。

\lecturefigure{slide-22-brain-unified-compute.jpg}{皮层同质性类比：统一学习算法的诱人猜想。}{Stanford Online 官方视频 00:14:42--00:15:09。}

\teachervoice{Karpathy 说视觉皮层、听觉皮层在宏观上相似，Transformer 在不同领域只改变部分“超参数”也许暗示一种统一强大学习算法。这是启发性类比，不是神经科学证据；架构迁移成功既可能来自普适计算结构，也可能来自现代数据与硬件共同偏好的工程解。}

\subsection{本章小结}

2011 年的 AI 被手工特征、领域词汇与专用 topology 分割；2012 年后的规模化深度学习统一了端到端优化语言；2017 年 Transformer 又把可复用主干推进到跨模态层面。下一章追问这套架构不是“从天而降”的：它沿着语言模型、seq2seq bottleneck 和 soft alignment 的问题链演化而来。

